%%%%%%%%%%%%%%%%%%%%%%%%%%%%%%%%%%%%%%%%%%%%%%%%%%%%%%%%%%%%%%%%%%%%%%%%%%%%%%%
%% Descr:       Vorlage für Berichte der DHBW-Karlsruhe
%% Author:      Prof. Dr. Jürgen Vollmer, vollmer@dhbw-karlsruhe.de
%% $Id: bericht.tex,v 1.19 2016/03/16 16:59:41 vollmer Exp $
%%  -*- coding: utf-8 -*-
%%%%%%%%%%%%%%%%%%%%%%%%%%%%%%%%%%%%%%%%%%%%%%%%%%%%%%%%%%%%%%%%%%%%%%%%%%%%%%%

\documentclass[
   ngerman          % neue deutsche Rechtschreibung
  ,a4paper          % Papiergrösse
% ,twoside          % Zweiseitiger Druck (rechts/links)
% ,10pt             % Schriftgrösse
% ,11pt
  ,12pt
  ,pdftex
%  ,disable         % Todo-Markierungen auschalten
]{report}

% Bitte die Codierung Ihrer Dateien auswählen:
% \usepackage[latin1]{inputenc}    % Für UNIX mit ISO-LATIN-codierten Dateien
% \usepackage[applemac]{inputenc}  % Für Apple Mac
  \usepackage[ansinew]{inputenc}   % Für Microsoft Windows
% \usepackage[utf8]{inputenc}        % UTF-8 codierte Dateien
                                   % Dieses Dokument ist unter Unix erstellt, daher
                                   % wird diese Input-Codierung benutzt.

\usepackage{bericht}
\usepackage{setspace}
\usepackage{listings}

%%%%%%%%%%%%%%%%%%%%%%%%%%%%%%%%%%%%%%%%%%%%%%%%%%%%%%%%%%%%%%%%%%%%%%%%%%%%%%%
%% Angaben zur Arbeit
%%%%%%%%%%%%%%%%%%%%%%%%%%%%%%%%%%%%%%%%%%%%%%%%%%%%%%%%%%%%%%%%%%%%%%%%%%%%%%%

\newcommand{\Autor}{Patrick Siewert}
\newcommand{\MatrikelNummer}{4363889}
\newcommand{\Kursbezeichnung}{TINF16B2}

\newcommand{\FirmenName}{Fiducia \& GAD IT AG}
\newcommand{\FirmenStadt}{Karlsruhe}
\newcommand{\FirmenLogoDeckblatt}{\fbox{\includegraphics[width=3cm]{lion}}}

% Falls es kein Firmenlogo gibt:
%  \newcommand{\FirmenLogoDeckblatt}{}

\newcommand{\BetreuerFirma}{Volker Werling}
\newcommand{\BetreuerDHBW}{Prof. Dr. Heinrich Braun}

%%%%%%%%%%%%%%%%%%%%%%%%%%%%%%%%%%%%%%%%%%%%%%%%%%%%%%%%%%%%%%%%%%%%%%%%%%%%%%%%%%%%%

\newcommand{\Was}{Projektarbeit}
% Wird auf dem Deckblatt in der Erklärung benutzt

%%%%%%%%%%%%%%%%%%%%%%%%%%%%%%%%%%%%%%%%%%%%%%%%%%%%%%%%%%%%%%%%%%%%%%%%%%%%%%%%%%%%%

\newcommand{\Titel}{Pr�sentation von PowerPoint-Folien durch einen Humanoiden
Roboter}
\newcommand{\AbgabeDatum}{3. April 2018}

\newcommand{\Dauer}{11 Wochen}

% \newcommand{\Abschluss}{Bachelor of Engineering}
\newcommand{\Abschluss}{Bachelor of Science}

% \newcommand{\Studiengang}{Informationstechnik}
\newcommand{\Studiengang}{Angewandte Informatik}

\hypersetup{%%
  pdfauthor={\Autor},
  pdftitle={\Titel},
  pdfsubject={\Was}
}

%%%%%%%%%%%%%%%%%%%%%%%%%%%%%%%%%%%%%%%%%%%%%%%%%%%%%%%%%%%%%%%%%%%%%%%%%%%%%%%

% Benutzt man das "biblatex"-Paket, dann muß das hier stehen:
% siehe auch die mit BIBLATEX markierten Zeilen in bericht.sty
\bibliography{bericht}

\begin{document}
\pagenumbering{roman}

%%%%%%%%%%%%%%%%%%%%%%%%%%%%%%%%%%%%%%%%%%%%%%%%%%%%%%%%%%%%%%%%%%%%%%%%%%%%%%%

\begin{titlepage}
\begin{center}
\vspace*{-2cm}
\FirmenLogoDeckblatt\hfill\includegraphics[width=4cm]{dhbw-logo}\\[2cm]
{\Huge \Titel}\\[2cm]
{\Huge\scshape \Was}\\[2cm]
{\large f�r die Pr�fung zum}\\[0.5cm]
{\Large \Abschluss}\\[0.5cm]
{\large des Studienganges \Studiengang}\\[0.5cm]
{\large an der}\\[0.5cm]
{\large Dualen Hochschule Baden-W�rttemberg Karlsruhe}\\[0.5cm]
{\large von}\\[0.5cm]
{\large\bfseries \Autor}\\[1cm]
{\large Abgabedatum \AbgabeDatum}
\vfill
\end{center}
\begin{tabular}{l@{\hspace{2cm}}l}
Bearbeitungszeitraum	         & \Dauer 			\\
Matrikelnummer	                 & \MatrikelNummer		\\
Kurs			         & \Kursbezeichnung		\\
Ausbildungsfirma	         & \FirmenName			\\
			         & \FirmenStadt			\\
Betreuer der Ausbildungsfirma	 & \BetreuerFirma		\\
Gutachter der Studienakademie	 & \BetreuerDHBW		\\
\end{tabular}
\end{titlepage}

%%%%%%%%%%%%%%%%%%%%%%%%%%%%%%%%%%%%%%%%%%%%%%%%%%%%%%%%%%%%%%%%%%%%%%%%%%%%%%%

% Nur f�r Bachelorarbeiten einf�gen:
%%%%%%%%%%%%%%%%%%%%%%%%%%%%%%%%%%%%%%%%%%%%%%%%%%%%%%%%%%%%%%%%%%%%%%%%%%%%%%%
%% Descr:       Vorlage für Berichte der DHBW-Karlsruhe, Erklärung
%% Author:      Prof. Dr. Jürgen Vollmer, vollmer@dhbw-karlsruhe.de
%% $Id: erklaerung.tex,v 1.6 2016/03/16 12:51:09 vollmer Exp $
%% -*- coding: utf-8 -*-
%%%%%%%%%%%%%%%%%%%%%%%%%%%%%%%%%%%%%%%%%%%%%%%%%%%%%%%%%%%%%%%%%%%%%%%%%%%%%%%

% In Bachelorarbeiten muss eine schriftliche Erklärung abgegeben werden.
% Hierin bestätigen die Studierenden, dass die Bachelorarbeit, etc.
% selbständig verfasst und sämtliche Quellen und Hilfsmittel angegeben sind. Diese Erklärung
% bildet das zweite Blatt der Arbeit. Der Text dieser Erklärung muss auf einer separaten Seite
% wie unten angegeben lauten.

\newpage
\thispagestyle{empty}
\begin{framed}
\begin{center}
\Large\bfseries Erklärung
\end{center}
\medskip
\noindent
Ich versichere hiermit, dass ich meine \Was\ mit
dem Thema: \enquote{\Titel} selbstständig verfasst und keine anderen als die
angegebenen Quellen und Hilfsmittel benutzt habe. Ich versichere zudem, dass die eingereichte elektronische Fassung mit der
gedruckten Fassung übereinstimmt.

\vspace{3cm}
\noindent
\underline{\hspace{4cm}}\hfill\underline{\hspace{6cm}}\\
Ort~~~~~Datum\hfill Unterschrift\hspace{4cm}
\end{framed}

\vfill
\emph{Sofern von der Ausbildungsstätte ein Sperrvermerk gewünscht wird, ist folgende Formulierung
zu verwenden:}
\begin{framed}
\begin{center}
\Large\bfseries Sperrvermerk
\end{center}
\medskip
\noindent
Der Inhalt dieser Arbeit darf weder als Ganzes noch in Auszügen Personen
auerhalb des Prüfungsprozesses und des Evaluationsverfahrens zugänglich gemacht
werden, sofern keine anders lautende Genehmigung der Ausbildungsstätte vorliegt.
\end{framed}

%%%%%%%%%%%%%%%%%%%%%%%%%%%%%%%%%%%%%%%%%%%%%%%%%%%%%%%%%%%%%%%%%%%%%%%%%%%%%%%
\endinput
%%%%%%%%%%%%%%%%%%%%%%%%%%%%%%%%%%%%%%%%%%%%%%%%%%%%%%%%%%%%%%%%%%%%%%%%%%%%%%%


\newpage
\tableofcontents           % Inhaltsverzeichnis hier ausgeben
\newpage
\listoffigures             % Liste der Abbildungen
\newpage
\listoftables              % Liste der Tabellen
\newpage
\lstlistoflistings         % Liste der Listings
\newpage
\listofequations           % Liste der Formeln

% Jetzt kommt der "eigentliche" Text
\newpage
%%%%%%%%%%%%%%%%%%%%%%%%%%%%%%%%%%%%%%%%%%%%%%%%%%%%%%%%%%%%%%%%%%%%%%%%%%%%%%
%% Descr:       Vorlage für Berichte der DHBW-Karlsruhe, Datei mit Abkürzungen
%% Author:      Prof. Dr. Jürgen Vollmer, vollmer@dhbw-karlsruhe.de
%% $Id: abk.tex,v 1.3 2016/03/16 12:21:40 vollmer draft $
%% -*- coding: utf-8 -*-
%%%%%%%%%%%%%%%%%%%%%%%%%%%%%%%%%%%%%%%%%%%%%%%%%%%%%%%%%%%%%%%%%%%%%%%%%%%%%%%

\chapter*{Abkürzungsverzeichnis}                   % chapter*{..} -->   keine Nummer, kein "Kapitel"
						         % Nicht ins Inhaltsverzeichnis
% \addcontentsline{toc}{chapter}{Akürzungsverzeichnis}   % Damit das doch ins Inhaltsverzeichnis kommt

% Hier werden die Abkürzungen definiert
\begin{acronym}[DHBW]
  % \acro{Name}{Darstellung der Abkürzung}{Langform der Abkürzung}
 \acro{Abk}[Abk.]{Abkürzung}

 % Folgendes benutzen, wenn der Plural einer Abk. benöigt wird
 % \newacroplural{Name}{Darstellung der Abkürzung}{Langform der Abkürzung}
 \newacroplural{Abk}[Abk-en]{Abkürzungen}

 \acro{H2O}[\ensuremath{H_2O}]{Di-Hydrog-Monoxid}

 % Wenn neicht benutzt, erscheint diese Abk. nicht in der Liste
 \acro{NUA}{Not Used Acronym}
 
 \acro{hri}[HRI]{Human-Robot Interaction}
 \acro{mqtt}[MQTT]{Message Queuing Telemetry Transport} 
\end{acronym}              % Abk�rzungsverzeichnis
\newpage
\onehalfspacing
\pagenumbering{arabic}
\chapter{Einleitung}\label{sec:einleitung}
\section{Die Fiducia \& GAD IT AG}\label{sec:fiducia-gad}
Die Fiducia \& GAD IT AG ist der Dienstleister f�r Informationstechnologie der
genossenschaftlichen FinanzGruppe. Das Unternehmen besch�ftigt aktuell rund
6.400 Mitarbeiter an den Verwaltungssitzen in Karlsruhe und M�nster und den
Gesch�ftsstellen in M�nchen, Frankfurt und Berlin. Die Fiducia \& GAD
erwirtschaftet einen j�hrlichen Konzernumsatz von rund 1,4 Milliarden Euro.

\subparagraph{}
Kunden der Fiducia \& GAD sind alle 1.000 Volksbanken und Raiffeisenbanken in
Deutschland und die Unternehmen der genossenschaftlichen FinanzGruppe. Au�erdem
geh�ren zahlreiche Privatbanken und Unternehmen anderer Branchen, wie z.B. der
ADAC, zum Kundenkreis der Fiducia \& GAD.

\subparagraph{}
Neben dem Betrieb der beiden Bankverfahren "`agree21"' und "`bank21"' in ihren
vier Hochsicherheitsrechenzentren, betreut die Fiducia \& GAD 173.000
Bankarbeitspl�tze und verwaltet knapp 83 Millionen Kundenkonten. Au�erdem stellt
die Fiducia \& GAD mit 36.000 eigenen Selbstbedienungsger�ten bundesweit eine
reibungslose Bargeldversorgung sicher. \cite{FiduciaGAD2018}

\section{Motivation}\label{sec:motivation}
Humanoide Roboter l�sen bei vielen Menschen eine gro�e Faszination aus. Obwohl
sie in Japan bereits weit verbreitet sind und mehrere Firmen immer
fortgeschrittenere entwickeln, trifft man im realen Leben nur sehr selten auf
humanoide Roboter. Die meisten Menschen kennen diese Art Roboter nur aus
Science-Fiction Filmen. Diese verbreiten ein faszinierendes, wenn auch teilweise
be�ngstigendes, Bild von Robotern, die den Menschen im t�glichen Leben
unterst�tzen und dabei in Bewegung, Sprache und Aussehen einem Menschen �hneln.
Doch humanoide Roboter sind nicht mehr nur Science-Fiction. Sie haben das
Potenzial t�gliche Begleiter der Menschen zu werden, wie zuletzt der PC, oder
als noch aktuellere technische Entwicklung, das Smartphone.

\subparagraph{}
Die Fiducia \& GAD erforscht M�glichkeiten humanoide Roboter produktiv
eizusetzen. Dazu sollen sie in eigene oder f�r ihre Kunden bereitgestellte
Gesch�ftsprozesse eingebunden werden. Zu diesen Zwecken wird der Roboter Pepper
von Softbank Robotics eingesetzt.

\section{Zielsetzung}\label{sec:zielsetzung}
Momentan werden hierbei zun�chst komplexere Anwendungsszenarien betrachtet,
beispielsweise eine gef�hrte Kontoer�ffnung. Diese werden im Rahmen von
Veranstaltungen immer auch wieder bei Kunden pr�sentiert. Hierbei ist meist die
Einbindung eigener, zur Veranstaltung passender Inhalte erw�nscht. Um dieser
Anforderung und der Usergruppe zu begegnen soll die M�glichkeit der
Pr�sentation von PowerPoint-Folien durch den Roboter geschaffen werden.

\subparagraph{}
Ziel ist es eine Online-Plattform zu erstellen, die PowerPoint-Dateien so
umwandelt, dass die Folien auf dem Tablet des Roboters angezeigt werden k�nnen.
Zus�tzlich soll Pepper die Notizen der einzelnen Folien jeweils zur
entsprechenden Folie vorlesen. Dazu wird ein Programm geschrieben, welches die
Folien zu Bildern und Text umwandelt. Au�erdem wird eine App entwickelt, die auf
dem Roboter installiert wird und die den Roboter die, von der Online-Plattform
umgewandelte und bereitgestellte Pr�sentation vortragen l�sst.

\newpage
\chapter{Grundlagen}\label{sec:grundlagen}
\section{Python}\label{sec:python}
"`Python ist eine portable, interpretative, objektorientiere
Programmiersprache"' \cite{Weigend2017}. Entwickelt wurde sie 1989 von Guido van
Rossum. Heute wird die Entwicklung von der nichtkommerziellen Organisation
"`Python Software Foundation"' \footnote{\url{https://www.python.org/psf/}}
koordiniert.

\subparagraph{}
Das Python-Skript wird von einem Interpreter ausgef�hrt. Python-Skripte k�nnen
auf verschiedenen Systemplattformen (Unix, Windows, Mac OS) laufen, deshalb
bezeichnet man Python als portable Sprache. Die aktuelle Version ist 3.6. Da auf
dem Roboter allerdings Python 2.7 installiert ist, wird die Anwendung mit dieser
Version entwickelt. Python ist kostenlos. \cite{Weigend2017}

\section{JavaScript}\label{sec:javascript}
JavaScript ist, wie der Name bereits sagt, eine Skriptsprache. Verwendet wird
sie haupts�chlich im Worl Wide Web. Mit JavaScript lassen sich Objekte, wie zum
Beispiel das Browserfenster, beeinflussen. \cite{Steyer2010}

\section{PowerPoint}\label{sec:powerpoint}


\section{Flask}\label{sec:flask}
Flask ist ein Python Package zum erstellen eines Servers.
\newpage
\chapter{Humanoide Roboter}\label{sec:humanoide-roboter}
\section{Allgemein}\label{sec:allgemein}
Erkl�rung was humanoide Roboter sind.

\section{Abgrenzung zu anderen Arten von Robotern}\label{sec:abgrenzung}
H�ufig ist der Unterschied zwischen humanoiden Robotern und anderen Arten von
Robotern nicht ganz klar. So verschwimmen teilweise die Grenzen zwischen
Industrierobotern, Servicerobotern, humanoiden Robotern und Androiden.

\subsection{Industrieroboter}\label{sec:industrieroboter}
Der Unterschied zwischen humanoiden Robotern und Industrierobotern ist relativ
deutlich. Industrieroboter werden dazu entwickelt einzelne Schritte oder den
gesamten Fertigungsprozess zu automatisieren. Ihre Bewegungsabl�ufe erinnern oft
an die eines Menschen, da Industrieroboter meist aus mehreren Gelenken bestehen,
die unabh�ngig voneinander bewegt werden k�nnen, �hnlich wie der menschliche
Arm. \cite{Weber2017} Bei Industrierobotern wird allerdings kein besonderer
Fokus auf menschliche Bewegungen gelegt. Im Gegensatz zu humanoiden Robotern
werden Industrieroboter nicht entwickelt um einem Menschen �hnlich zu sein,
sondern um m�glichst effizient arbeiten zu k�nnen.

\subsection{Serviceroboter}\label{sec:serviceroboter}
Unter Servicerobotern versteht man Roboter, die eine Aufgabe gr��ten Teils
autonom erledigen. Bei diesen Aufgaben handelt es sich um Dinge, die ein Mensch
nicht erledigen kann oder will. Weit verbreitete Serviceroboter sind z.~B.
Staubsaugroboter oder M�hroboter. Der Unterschied zu humanoiden Robotern besteht
hierbei darin, dass Serviceroboter nur eine Art von Aufgaben erledigen k�nnen.
Au�erdem ist ihr Aussehen an die zu erledigende Aufgabe angepasst, nicht an das
Aussehen eines Menschen.

\subsection{Androiden}\label{sec:androiden}
Bei Androiden handelt es sich um humanoide Roboter, deren Zweck es ist, Menschen
m�glichst gut nachzuahmen. Wenn �ber das Thema Roboter im Allgemeinen gesprochen wird,
denken viele Menschen zuerst an Androiden oder humanoide Roboter auch wenn heute
Industrieroboter viel weiter verbreitet sind und die Chance z.~B. einen
M�hroboter zu besitzen deutlich h�her ist, als einen Androiden zu Hause zu
haben. Dass Androiden trotzdem so verbreitet in der Vorstellung der Menschen
sind, liegt zum Einen an B�chern, Serien und Filmen und zum Anderen an der
Faszination, die Androiden ausl�sen. \cite{Dautenhahn2011}

\subsubsection{Androiden in Medien}
In vielen Science-Fiction B�chern, Serien und Filmen, sind Androiden oder
humanoide Roboter ein wesentlicher Bestandteil der Geschichte. Dabei treten sie
in sehr unterschiedlichen Arten auf. Jedoch haben alle gemeinsam, dass sie in
nicht all zu ferner Zukunft existieren und einen gewissen Einfluss auf das Leben
der Menschen haben. Dieser Einfluss ist in manchen Darstellungen positiv, in
anderen negativ. Mehr hierzu in Kapitel \ref{sec:roboter-in-medien} auf Seite
\pageref{sec:roboter-in-medien}.

\subparagraph{}
Geht es nach den Vorstellungen der Autoren,
scheint ein Leben ohne humanoide Roboter in Zukunft nicht vorstellbar zu sein.
Im realen Leben ist die Wahrscheinlichkeit mit einem Industrieroboter oder einem
Serviceroboter in Kontakt zu kommen zwar deutlich h�her als auf einen Androiden
zu treffen, durch das h�ufige Auftreten in den Medien sind Androiden jedoch
sehr weit in der Vorstellung der Menschen vertreten.

\subsubsection{Faszination von Androiden}
Obwohl Serviceroboter oder Industrieroboter technisch sehr aufw�ndig sein
k�nnen, haben sie f�r einen einzelnen Menschen meist keine sehr gro�e Bedeutung.
Ein M�hroboter ist zwar n�tzlich, jedoch tut er von au�en betrachtet nichts
anderes als den ganzen Tag �ber den Rasen zu fahren. Mit Industrierobotern
kommen Menschen noch seltener in Kontakt, wenn sie nicht gerade in einer Firma
arbeiten, in der diese eingesetzt werden. Auch diese Industrieroboter erledigen
meist nur eine einzelne Aufgabe und sind f�r Ausenstehende auf dauer wenig
interessant.

\subparagraph{}
Androiden l�sen hier eine viele gr��ere Faszination aus. Sie
stellen eine Zukunftsvision dar, die in der Realit�t noch kaum vertreten ist.
Viele Menschen sind neugierig, was Androiden alles k�nnen und sind fasziniert
von ihren F�higkeiten. Au�erdem k�nnen Menschen mit Androiden interagieren.
Androiden k�nnen auf Fragen antworten, auf Ber�hrungen reagieren und die
Menschen unterhalten. Dadurch sind Androiden auf Dauer viel spannender als nicht
interaktive Industrie- oder Serviceroboter.

\subsubsection{Androiden in der Forschung}
Trotz der weiten Verbreitung in Medien und der Vorstellung von
Menschen, machen Androiden nur einen kleinen Teil der Forschung innerhalb der
viel gr��eren Bereiche Robotik und K�nstliche Intelligenz aus.
\cite{Dautenhahn2011} Trotzdem arbeiten verschiedene Unternehmen an der
Entwicklung von humanoiden Robotern und Androiden. Beipiele hierf�r sind Pepper
(Kapitel \ref{sec:pepper}) und Nao von SoftbankRobotics und Aldebaran oder Asimo
von Honda, welcher als der am weitesten entwickelte humanoide Roboter der Welt
beworben wird. \cite{Honda2018}

\section{Moderne Humanoide Roboter}
\subsection{Pepper}\label{sec:pepper}
Pepper ist ein humanoider Roboter der Firmen SoftBank Robotics und Aldebaran.
Vom Hersteller wird er als freundlich, liebenswert und �berraschend beworben.
Entwickelt wurde Pepper um ein allt�glicher Begleiter f�r Menschen zu sein. Er
wird als "`viel mehr als ein Roboter"' beschrieben. Mit seiner K�rpersprache und
der Stimme soll er auf die "`nat�rlichste und intuitivste"' Art mit Menschen
kommunizieren. \cite{SoftBank2018}

\subparagraph{}
Mit einer Gr��e von etwa 1,20m, seinem kindlichen Gesicht und beweglichen Armen
und H�nden, soll Pepper Menschen gl�cklich machen. Sein Kopf bewegt sich, um
vorbeilaufende Passanten anzuschauen oder einem Menschen in die Augen zu
schauen, wenn er mit dem Roboter spricht. \cite{Markowitz2015}

\subparagraph{}
Das Design von Pepper ist darauf ausgelegt Emotionen sowohl zu verstehen als
auch auszul�sen. Die dazu verwendeten Eigenschaften sind auch wichtig f�r das
Vorstellen von PowerPoint Pr�sentation.

\subsubsection{H�ren und Sprechen}\label{sec:hoeren-und-sehen}
Mit Lautsprechern und Mikrofonen ausgestattet, kann Pepper gesprochenes
Verstehen und darauf reagieren. In der Anwendung soll Pepper die Notizen einer
PowerPoint Pr�sentation vorlesen, wozu die Mikrofone benutzt werden. Die
F�higkeit zuzuh�ren soll dazu verwendet werden Pepper w�hrend der Pr�sentation
Anweisungen geben k�nnen. Er soll auf Befehle wie "`Pause"' und "`Weiter"'
reagieren und die Pr�sentation entsprechen pausieren bzw. fortsetzen.

\subsubsection{Sehen}\label{sec:sehen}
Eine 3D-Kamera und zwei HD-Kameras helfen Pepper seine Umgebung zu analysieren
und Menschen und Bewegungen zu erkennen. W�hrend einer Pr�sentation wird diese
F�higkeit, au�er bei der schon vom Hersteller implementierten
Kollisionserkennung, nicht verwendet. Es wird nicht n�tig sein spezielle Dinge
zu erkennen.

\subsubsection{Verbindung}\label{sec:verbindung}
Pepper verf�gt �ber eine Internetverbindung. Dies ist wichtig, da die PowerPoint
Pr�sentation vor der Vorstellung durch Pepper nicht zuerst auf den Roboter
geladen werden soll. Die Informationen wie die einzelnen Folien und der
zugeh�rige Text sollen auf einem Server bereitgestellt werden auf den der
Roboter zugreifen kann. Mehr hierzu in Kapitel \ref{sec:umsetzung} auf Seite
\pageref{sec:umsetzung}.

\subsubsection{Tablet}\label{sec:tablet}
Auf der Brust von Pepper ist ein Tablet angebracht, auf dem Bilder, Videos oder
Webseiten angezeigt werden k�nnen. Auf diesem Tablet sollen die einzelnen Folien
der PowerPoint Pr�sentation angezeigt werden. \cite{SoftBankII2018}

\section{Human-Robot Interaction}\label{sec:hri}
\ac{hri} ist ein relativ neues, aber wachsendes Forschungsgebiet, welches sich
mit der Interaktion zwischen Menschen und Robotern befasst. Au�erdem soll
herausgefunden werden wie Roboter am besten mit Menschen zusammenarbeiten
k�nnen. Dabei hat \ac{hri} nicht nur Einfluss auf die Wirtschaft, sondern auch
auf m�gliche Arten der Beziehungen zwischen Menschen und Robotern. Deshalb
verbindet \ac{hri} verschiedene Wissenschaften, wie Psychologie und
Sozialwissenschaften mit Informatik und Robotik. Eines der Hauptziele von
\ac{hri} ist das Erforschen m�glichst nat�rlicher Wege der Kommunikation
zwischen Menschen und Robotern. \cite{Dautenhahn2011}

\subsection{Nat�rliche Formen der
Kommunikation}\label{sec:nat-formen-kommunikation}
Um eine nat�rliche Kommunikation zwischen Menschen und Robotern zu erm�glichen,
m�ssen mehrere, von einem Benutzer potenziell verwendbare, Interfaces
bereitgestellt werden. Dabei werden klassische mit neueren Interfaces
kombiniert.

\subsubsection{Graphisches Interface}
Ein klassisches Interface zur Interaktion mit Maschinen sind graphische
Input-Output Schnittstellen. Wie bereits von der Kommunikation mit Computern
bekannt, k�nnte ein Roboter durch Eingaben mit Maus und Tastatur gesteuert
werden. Diese Form des Interfaces ist durch die weite Verbreitung von Computern
den meisten Menschen bekannt. Allerdings sind Maus und Tastatur keine
nat�rliche Art der Kommunikation. Dieses Interface kann durch die Verwendung
eines Touchscreens intuitiver und damit auch nat�rlicher gestaltet werden.

\subparagraph{}
Vorteil von graphischen Interfaces ist, dass der Benutzer dem Roboter klare
Instruktionen geben kann. �ber Men�s muss der Benutzer eine bestimmte Aktion
ausw�hlen, die der Roboter ausf�hren soll. Der Roboter f�hrt dann den
entsprechenden Programmcode aus. Auf diese Weise gibt es keinen Spielraum f�r
eventuelle Fehlinterpretationen der Instruktionen durch den Roboter.

\subsubsection{Sprache}
Sprache ist eine sehr nat�rliche Form der Kommunikation. Ein Mensch kann einem
Roboter einen Befehl geben, indem er dem Roboter sagt, welche Aktion er
ausf�hren soll. Der Mensch ist es gewohnt durch das Sprechen zu kommunizieren,
deshalb ist Sprache sowohl ein nat�rliches als auch ein intuitives Interface.

\subparagraph{}
Allerdings ist die Umsetzung dieses Interface komplizierter als die eines
einfachen graphischen Interface. Zun�chst muss der Roboter die Worte, die der
Mensch spricht erkennen. Spracherkennung ist bereits weit verbreitet und wird in
verschiedenen Anwendungen verwendet. So lassen sich z.~B. Smartphones mit Siri
oder dem Google Assistant per Sprachbefehlen bedienen. Allerdings muss der
Roboter auch das verstehen, was der Mensch meint. Schon eine vermeintlich
einfache Anweisung, wie z.~B. "`Setz dich."' kann vom Roboter unterschiedlich
interpretiert werden. Er kann sich z.~B. entweder auf einen Stuhl oder auf den
Boden setzen.

\subparagraph{}
Zwar sind Sprachbefehle nat�rlicher und intuitiver als graphische Eingaben auf
einem Display, dem Menschen, der dem Roboter Anweisungen gibt, wird es jedoch
schwerer fallen seine Anweisungen so pr�zise zu formulieren, dass der Roboter
genau das tut, was er tun soll. Durch die Auswahl aus einem Men� w�re dies
einfacher.

\subparagraph{}
Au�erdem muss der Roboter sich an die Dynamik, die sich aus einer Interaktion
ergibt anpassen. Kommunikation besteht meist nicht nur aus einem Befehl des
Menschen und der darauf folgenden Aktion des Roboters. Der Roboter kann
nachfragen, wenn ein Befehl nicht pr�zise genug formuliert wird und entsprechend
auf die Antwort des Menschen reagieren. Diese Frage-Antwort Dynamik muss der
Roboter verstehen und Antworten in den richtigen Kontext setzen.

\subsubsection{Visuelle Kommunikation}
Mit einer Kamera kann der Roboter Bewegungen, Gesten und Mimik eines Menschen
erkennen. So kann der Roboter mit der Hand gesteuert werden, reagieren, wenn ein
Mensch Blickkontakt zu ihm aufnimmt. So ist es auch m�glich einem Roboter einen
Bewegungsablauf vorzumachen, den er dann wiederholen kann, was das Erkl�ren
einer Aktion deutlich vereinfachen kann.

\subsubsection{Tastsensoren}
In der direkten Interaktion wird es auch zu
Ber�hrungen zwischen Meschen und Robotern oder Robotern untereinander kommen. Um
dabei keine Sch�den zu verursachen oder Menschen zu verletzen muss der Roboter
auf Ber�hrungen reagieren k�nnen. Ber�hrungen sind eine direkte Form der
Kommunikation, die dem Roboter unmissverst�ndlich bestimmte Anweisungen geben
kann. So kann fest einprogrammiert werden dass sich der Roboter nicht mehr
bewegt, nachdem er ber�hrt wurde. Damit wird verhindert, dass er Menschen
verletzt. \cite{Prassler2004}

\subsection{Roboter in den Medien}\label{sec:roboter-in-medien}
Die Meisten Menschen bekommen nur durch Medien eine Vorstellung von humanoiden
Robotern, da sie in der Realit�t kaum vertreten sind. Au�er bei besonderen
Veranstaltungen ist die Chance auf einen humanoiden Roboter zu treffen sehr
gering. Am h�ufigsten treten humanoide Roboter im Genre Science-Fiction auf.
Dabei kann zwischen dystopischen Geschichten, in denen Roboter die Menschheit
unterdr�cken oder vernichten, und Geschichten mit "`guten"' Robotern
unterschieden werden.

\subsubsection{B�se Roboter}

??? Soll ich �ber so etwas schreiben, oder geht das zu sehr in Richtung
Philosophie, bzw. passt das in eine Praxisarbeit in Informatik ???

\subsubsection{Gute Roboter}

??? Soll ich �ber so etwas schreiben, oder geht das zu sehr in Richtung
Philosophie, bzw. passt das in eine Praxisarbeit in Informatik ???
\newpage
\chapter{Umsetzung}\label{sec:umsetzung}

\section{Struktur}\label{sec:struktur}
Der gr��te Teil der Anwendung l�uf auf einem Flask-Server. Auf diesem Server
werden PowerPoint-Dateien, die vom Benutzer hochgeladen werden k�nnen, zur
Pr�sentation durch Pepper umgewandelt. Dazu werden die einzelnen Folien zu
Bildern zur Darstellung als HTML-Datei konvertiert. Die Notizen der Folien
werden als Text extrahiert. Diesen Text wird Pepper w�hrend der Pr�sentation
passend zur jeweiligen Folie vorlesen. Die Folien werden auf dem Tablet auf
Peppers Brust angezeigt. Welche Pr�sentation vorgestellt werden soll, kann auf
dem Tablet ausgew�hlt werden.

\begin{lstlisting}[float, language=Python, frame=single, framexleftmargin=15pt,
style=algoBericht, label={lst:vorlage}, captionpos=b, caption={Vorlage f�r das
Einf�gen von Code-Beispielen}]
print "Hello, world!"
\end{lstlisting}

\section{Server}\label{sec:server}

\newpage
\chapter{Fazit}\label{sec:fazit}
Hier steht ein Fazit
\newpage
\chapter{Ausblick}\label{sec:ausblick}
Hier steht ein Ausblick
\newpage

% Ab hier beginnt der Anhang
\appendix
\addcontentsline{toc}{chapter}{Anhang}

\addcontentsline{toc}{chapter}{Index}
\printindex

\addcontentsline{toc}{chapter}{Literaturverzeichnis}

% Haben Sie das "biblatex"-Paket nicht installiert, benutzen Sie folgendes:
% Ohne das "biblatex"-Paket (s. bericht.sty) produziert folgendes
% "deutsche" Zitate in Literaturverzeichnissen gemaß der Norm DIN 1505,
% Teil 2 vom Jan. 1984.
% Die Zitatmarken werden alphabetisch nach Verfassern
% sortiert und sind durch abgekürzte Verfasserbuchstaben plus
% Erscheinungsjahr in eckigen Klammern gekennzeichnet.

% \bibliographystyle{alphadin}
% \bibliography{bericht}

%%%%%%%%%%%%%%%%%%%%%%%%%%%%%%%%%%%%%%%5
% BIBLATEX
% Benutzt man das "biblatex"-Paket, muß man folgendes schreiben:
\def\refname{Literaturverzeichnis}
\printbibliography
%%%%%%%%%%%%%%%%%%%%%%%%%%%%%%%%%%%%%%%5
\newpage
\addcontentsline{toc}{chapter}{Liste der ToDo's}
\listoftodos[Liste der ToDo's]


\end{document}